\bookchapter{Глава 6: Защита вашей цифровой жизни — Безопасность и конфиденциальность}

\chapterintro{Безопасность телефона = конфиденциальность + сохранность средств. В эпоху, когда данные — это актив, защита телефона означает защиту себя.}

\booksection{Блокировка экрана: первая физическая линия обороны}

Блокировка экрана — первый рубеж защиты вашего телефона, и при этом самый недооценённый.

\booksubsection{Рейтинг уровня безопасности:}

\begin{enumerate}
  \item \textbf{Пароль / PIN-код}: наивысшая безопасность, последняя гарантия (на него полагаются, когда забыт отпечаток/лицо).
  \item \textbf{Сканер отпечатков пальцев}: идеальный баланс безопасности и удобства.
  \item \textbf{Распознавание лица}: самое удобное, но безопасность зависит от технологии (3D-структурный свет безопасен, 2D-фото легко обмануть).
  \item \textbf{Графический ключ}: удобно, но легко оставляет следы, по которым можно угадать.
\end{enumerate}

\booksubsection{Настройка надёжной блокировки:}

\begin{itemize}
  \item Обязательно установите сложный пароль (не дату рождения и не 123456).
  \item Включите биометрию.
  \item Время автоблокировки установите на «Сразу» или минимальное.
\end{itemize}

\booksection{«Найти мой телефон»: главный инструмент защиты от кражи}

\textbf{Обязательно включите!}Он позволяет:

\begin{itemize}
  \item \textbf{Определить местоположение}: где потерян телефон — найти на карте.
  \item \textbf{Воспроизвести звук}: заставить телефон звонить — удобно для поиска дома.
  \item \textbf{Режим пропажи}: удалённо заблокировать телефон и показать контактный номер.
  \item \textbf{Удалённое стирание}: крайняя мера — удалить все данные (но после стирания отслеживание невозможно).
\end{itemize}

Зайдите в [Настройки], на iPhone ищите «Найти iPhone», на Android — «Найти устройство» или аналогичный сервис бренда, убедитесь, что он включён (зелёный индикатор).

\booksection{Обновление системы: ставим заплатки, закрываем дыры}

Главная цель обновлений системы — исправление уязвимостей безопасности, как заделка дыр в стене дома. Поэтому, увидев уведомление об обновлении системы, по возможности обновляйтесь сразу.

Конечно, крупные обновления версий (например, с iOS 17 на iOS 18) могут иметь проблемы совместимости — можно подождать несколько дней перед обновлением.

\booksection{Распознавание и предотвращение телекоммуникационного мошенничества}

В эпоху стремительного развития AI-технологий традиционная модель «верить тому, что видишь» полностью утратила эффективность. Мошенники используют дипфейки (Deepfake) и другие AI-технологии для перехода от массового мошенничества к точечной охоте на жертв.

\booksubsection{Самые распространённые схемы мошенничества в эпоху AI}

\textbf{(1) Мошенничество с подменой лица и голоса через AI}\par

Злоумышленники незаконно собирают фотографии высокого разрешения из соцсетей жертвы или всего несколько секунд голосового образца и с помощью технологии дипфейков создают неотличимые от реальных «видео с родственниками» или «голосовые сообщения». Типичные сценарии:\textbf{внезапная авария, срочный заём, похищение}и т.д. Такое мошенничество использует эмоциональные связи — семейные и дружеские — и\textbf{крайне трудно распознать}.

\textbf{(2) Мошенничество под видом официальных лиц / службы поддержки}\par

AI-технологии способны точно имитировать голос и интонации официальных служб поддержки и сотрудников правоохранительных органов, а также синтезировать крайне реалистичные поддельные скриншоты переводов, электронные документы, ордера и т.д.,\textbf{создавая панику или давление авторитетом}, чтобы заставить жертву перевести деньги или включить демонстрацию экрана, похищая таким образом данные учётной записи.

\textbf{(3) AI-мошенничество с инвестициями}\par

С помощью AI-алгоритмов создаются поддельные скриншоты «гарантированной прибыли», интерфейсы виртуальных программ для подбора акций или «экспертные» инвестиционные отчёты. Под предлогом высокой доходности продаются дорогие инвестиционные услуги или жертву склоняют перевести средства на контролируемый мошенниками торговый счёт —\textbf{в итоге деньги исчезают}.

\booksubsection{Методы разоблачения — принцип «волшебного зеркала»}

AI может клонировать голос и лицо, но он не безупречен, есть несколько уязвимых мест:

\begin{itemize}
  \item \textbf{Метод проверки движением}: при видеозвонке активно попросите собеседника выполнить заранее оговорённую последовательность движений, например: «проведи рукой перед лицом», «покажи профиль с ухом», «быстро моргни 5 раз подряд». Современные технологии дипфейков\textbf{крайне подвержены зависаниям, искажениям лица и артефактам на границах}при обработке сложных динамических перекрытий, быстрых движений и изменений освещения.
  \item \textbf{Метод секретного кодового слова}: установите и используйте заранее оговорённые приватные проверочные вопросы, например: «Что ты заказал в прошлый раз, когда мы вместе ужинали?», «Какой был мой первый подарок тебе?» — детали, известные только вам двоим. Если собеседник заметно медлит или отвечает невпопад —\textbf{немедленно завершите звонок и запустите протокол безопасности}.
  \item \textbf{Метод двойной проверки}: строго соблюдайте принцип\textbf{«Четыре НЕ и одна проверка»}, опубликованный Управлением кибербезопасности Министерства общественной безопасности: не доверять незнакомой информации, не переводить деньги бездумно, не разглашать личные чувствительные данные, не распространять подозрительный контент. При любых денежных операциях обязательно\textbf{перезвоните по известному оригинальному номеру или подтвердите лично}. При переводе рекомендуется выбирать режим «зачисление через 24 часа», чтобы оставить запас времени для выявления мошенничества и возврата средств.
\end{itemize}

\booksubsection{Традиционные распространённые схемы (всё ещё актуальны):}

\begin{itemize}
  \item \textbf{Выигрыш / Кэшбэк за накрутку}: заманивают выгодой, сначала дают немного, потом обманывают на крупную сумму.
  \item \textbf{Мошенничество в онлайн-знакомствах}: «разделка свиньи» — входят в доверие, а потом выманивают деньги.
  \item \textbf{Выдача себя за знакомого}: взлом аккаунта и просьба денег в долг (в эпоху AI стало ещё сложнее отличить).
\end{itemize}

\booksubsection{Встроенная защита телефона:}

\begin{itemize}
  \item Блокировка спам-звонков.
  \item Защита от кражи проверочных кодов.
  \item Проверка безопасности ссылок в SMS.
\end{itemize}

\booksubsection{Мантра защиты от мошенничества:}

\textbf{Не доверять, не разглашать, не нажимать, не переводить, сомневаешься — сначала проверь!}Столкнувшись со «срочным делом», сначала успокойтесь на 5 минут — почти любое мошенничество раскроется в этот момент.

\booksection{Защита от «доксинга» и атак социальной инженерии: разрываем цепь утечки личных данных}

Суть «доксинга» (атаки социальной инженерии) заключается в\textbf{систематическом сборе, анализе и сопоставлении цифровых фрагментов личности для построения полного профиля}. Ключ к защите — провести чёткую красную линию для личной конфиденциальности и выстроить многоуровневую систему обороны.

\booksubsection{Очистка метаданных фотографий}

Базовые методы очистки EXIF описаны в Главе 2 «Управление изображениями и файловой системой». Перед публикацией фото в соцсетях необходимо дополнительно\textbf{полностью удалить метаданные фотографии (информацию EXIF)}, включая время съёмки, модель устройства, координаты GPS и другие чувствительные данные. Также необходимо тщательно замазать или размыть\textbf{номера домов, автомобильные номера, штрихкоды билетов, номера удостоверений личности}и другую ключевую информацию на фото, чтобы случайно не раскрыть маршруты передвижения.

\begin{itemize}
  \item iOS: [Настройки] > [Конфиденциальность и безопасность] > [Камера] — можно отключить «Точное местоположение»; при отправке выберите «Только текущее фото» для удаления метаданных.
  \item Android: некоторые приложения галереи (например, Google Фото) при отправке спрашивают, сохранять ли информацию о местоположении.
  \item Универсальные инструменты: ExifTool, Metapho и др. для просмотра/редактирования EXIF.
\end{itemize}

\booksubsection{Ограничение разрешений приложений}

Базовые методы управления разрешениями описаны в Главе 4 «Управление разрешениями». В контексте защиты от доксинга необходимо дополнительно регулярно проверять и отключать ненужные разрешения у приложений электронной коммерции и утилит —\textbf{микрофон, контакты, местоположение}— чтобы приложения не «подслушивали» повседневные разговоры в фоне и незаконно не собирали сеть личных контактов; давно не используемые приложения с высокими разрешениями следует своевременно удалять, сокращая источники утечки конфиденциальности.

\begin{itemize}
  \item При запросе «Доступ к контактам» сначала спросите: действительно ли этому приложению нужно знать всех моих контактов?
  \item Приложения, которыми вы не пользовались три месяца, решительно удаляйте. Разрешения со временем должны ужесточаться, а не ослабляться.
\end{itemize}

\booksubsection{Защита умных устройств}

Умные камеры, умные замки, умные колонки и другие IoT-устройства при первом использовании требуют\textbf{немедленной смены пароля по умолчанию}— установите надёжный уникальный пароль с буквами в верхнем и нижнем регистре, цифрами и спецсимволами; редко используемые функции удалённого доступа следует держать отключёнными; регулярно проверяйте обновления прошивки и своевременно закрывайте уязвимости безопасности,\textbf{чтобы не дать хакерам проникнуть и шпионить за частным пространством дома}.

\begin{tcbwarn}{Реальный случай-предупреждение}
В 2022 году полиция одного из регионов сообщила, что хакеры, взломав пароли по умолчанию домашних камер определённого бренда (многие пользователи никогда их не меняли), похитили изображения из спален тысяч домохозяйств и продавали их в интернете.\textbf{Не менять пароль по умолчанию = оставить дверь дома открытой}.

\end{tcbwarn}

\booksection{Защита конфиденциальности: ваши данные — ваша собственность}

Фотографии, история переписки, треки местоположения в телефоне — по сути, это ваши «цифровые активы». Утечка грозит не только рекламным спамом, но и доксингом, мошенничеством и даже шантажом.

\booksubsection{Принцип минимизации разрешений приложений:}

\begin{itemize}
  \item \textbf{Местоположение}: для такси/карт — «Всегда»; для остальных приложений — «Только при использовании».
  \item \textbf{Контакты}: давать, только если необходимо (например, WeChat для добавления друзей), иначе всегда запрещать.
  \item \textbf{Фото / Все фото}: начиная с iOS 14+ можно разрешить доступ к «Выбранным фото», Android 13+ также поддерживает точную авторизацию.
  \item \textbf{Микрофон / Камера}: обычно отключены, включать по необходимости. При подозрении на скрытую запись смахните в центр управления и проверьте наличие «зелёной/оранжевой точки».
  \item \textbf{Фоновое отслеживание}: начиная с iOS 14.5+ приложения обязаны запрашивать разрешение перед использованием IDFA; на Android можно отключить «Рекламный ID».
\end{itemize}

\booksubsection{Практика управления конфиденциальностью:}

\begin{itemize}
  \item \textbf{Отчёт о конфиденциальности (iOS)}: [Настройки] > [Конфиденциальность и безопасность] > [Отчёт о конфиденциальности приложений] — посмотрите, какие приложения тайно обращаются к вашему местоположению, камере, микрофону.
  \item \textbf{Журнал использования разрешений (Android)}: [Настройки] > [Конфиденциальность] > [Журнал использования разрешений] — посмотрите, кто и когда вызывал чувствительные разрешения.
  \item \textbf{Ограничение рекламного отслеживания}: отключите в системных настройках «Рекламный ID» или «Ограничить рекламное отслеживание» — это уменьшит навязчивую таргетированную рекламу.
\end{itemize}

\booksection{Чек-лист «Скорой помощи» при утере / краже}

Если телефон потерян — не паникуйте, действуйте в следующем порядке:

\begin{enumerate}
  \item \textbf{Немедленно заблокируйте SIM-карту}: свяжитесь с оператором для блокировки, чтобы её не использовали для получения проверочных кодов.
  \item \textbf{Удалённая блокировка}: через облачный сервис (iCloud / облако бренда) включите «Режим пропажи».
  \item \textbf{Смените важные пароли}: немедленно смените пароли от всех важных аккаунтов, в которые вы входили с телефона (платежи, банк, почта).
  \item \textbf{Заявление в полицию}: подайте заявление в отделение полиции, запишите номер IMEI. Это также поможет подтвердить личность при нахождении телефона.
  \item \textbf{Оповестите родных и друзей}: чтобы мошенники не просили денег от вашего имени.
  \item \textbf{(В последнюю очередь) Удалённое стирание}: выполняйте только когда уверены, что телефон не найти — после стирания функция отслеживания отключается.
\end{enumerate}

\booksection{Безопасность SIM-карты и eSIM}

\begin{itemize}
  \item \textbf{Установите PIN-код SIM-карты}: после включения при каждой перезагрузке или смене карты потребуется ввод PIN-кода, что предотвращает обход блокировки экрана извлечением карты.
  \item \textbf{Преимущества eSIM}: нет физической карты, которую можно извлечь, — безопасность выше; одновременно поддерживает переключение между несколькими номерами, очень удобно в путешествиях.
  \item \textbf{Риск вторичной замены карты}: остерегайтесь мошенничества с «вторичной заменой SIM-карты оператором» — злоумышленники выдают себя за вас и перевыпускают SIM-карту, захватывая ваш номер. Самая эффективная защита — установить требование «перевыпуск карты только при личном присутствии владельца с удостоверением личности в офисе оператора».
\end{itemize}

\begin{tcbwarn}{Чего категорически нельзя делать}
\begin{itemize}
  \item Использовать слабые пароли вроде даты рождения / 123456 / qwerty;
  \item Отправлять проверочный код кому бы то ни было (даже «службе поддержки»);
  \item Нажимать на подозрительные ссылки в SMS, особенно от зарубежных номеров, начинающихся на 00;
  \item Входить в интернет-банк или платёжные системы через публичный Wi-Fi;
  \item Предоставлять незнакомцам «демонстрацию экрана» или «удалённое управление».
\end{itemize}

\end{tcbwarn}

\begin{tcbpractice}{Попробуйте включить и просимулировать «удалённую блокировку»}
\begin{enumerate}
  \item Убедитесь, что функция «Найти мой телефон» включена.
  \item В браузере компьютера или на другом телефоне войдите в соответствующий облачный аккаунт (iCloud.com или сайт облачного сервиса бренда).
  \item Найдите функцию «Найти устройство» и попробуйте запустить для своего телефона команду «Воспроизвести звук» или «Заблокировать», чтобы ознакомиться со всем процессом.
\end{enumerate}

\end{tcbpractice}
